\documentclass[11pt]{article}
\usepackage[a4paper,margin=27mm]{geometry}
\usepackage{amsmath,amssymb,amsthm,hyperref}
\newtheorem{theorem}{Theorem}
\newtheorem{lemma}{Lemma}
\newcommand{\E}{\mathbb E}
\newcommand{\Pol}{\operatorname{Pol}}
\title{A power lower bound for every permutation-fair die}
\author{Research proof: independently reviewed}
\date{30 September 2026}
\begin{document}\maketitle
\noindent\textbf{Status.} This version proves the independently reviewed $7/6$ exponent.
Its essential input is the finite mean-deficit balance
theorem in
\path{rational_designs/asymptotic/finite_endpoint_odds_balance.tex}.

\begin{theorem}\label{main}
There is an absolute constant $c>0$ such that every individual die in
a permutation-fair collection of $n$ equiprobable finite dice has at
least $c n^{7/6}$ faces.
More generally, suppose $p_1,\ldots,p_m:[K]\to[0,1]$ are nondecreasing,
the rows have positive probability weights $w_q$, and
\[
 \E_w\prod_{j\in S}p_j=\frac1{|S|+1}\qquad(S\subseteq[m]).
\]
Then $w_1,w_K\le C m^{-7/6}$ for an absolute $C$.
\end{theorem}

All constants denoted by $C,c$ below are absolute and may change from
line to line. Fixed numerical choices, such as $\theta=1/16$, do not
depend on $m$ or on the representation. We prove the assertion for
sufficiently large $m$; increasing $C$ covers the remaining range.

\section*{The finite balance input}
Put $u_j=1-p_j$ and $s=\sum_j u_j$. The cited finite balance theorem
and its row-mass corollary give, for $m\ge256$,
\begin{align}
 \max_j u_j&\le256(s+1)/m,                    \label{balance}\\
 \Pr_w(s\le S)&\le512(S+1)/m\quad(S\ge0).          \label{mass}
\end{align}
The first bound is the cited theorem when $s\le m/256$; otherwise
it follows simply from $u_j\le1$. In particular, for every fixed $a>0$,
\begin{equation}\label{laplace}
 \E_w e^{-as}\le C_a/m.
\end{equation}
We shall also use the following uniform consequence for integers $k\ge0$:
\begin{equation}\label{weightedmoments}
 \E_w (s+1)^k e^{-as}\le \frac{C_a^{k+1}k!}{m}.
\end{equation}
Indeed split $e^{-as}=e^{-as/2}e^{-as/2}$, apply \eqref{laplace}, and
use $\sup_{s\ge0}(s+1)^k e^{-as/2}\le C_a^k k!$.

\section*{A stable polarization comparison}
For a polynomial $P(z)=\sum_k a_kz^k$ of degree at most $N$, write
\[
 \Pol_N(P)(z_1,\ldots,z_N)
    =\sum_k a_k\frac{e_k(z_1,\ldots,z_N)}{\binom Nk}.
\]
This is its symmetric multiaffine polarization.

\begin{lemma}[Centered elementary coefficients]\label{elementary}
If $\sum_{i=1}^Ny_i=0$ and $|y_i|\le R$, then for $2\le k\le N/2$,
\[
 \frac{|e_k(y)|}{\binom Nk}
       \le \left(\sqrt{2e}\,R\sqrt{k/N}\right)^k.
\]
\end{lemma}
\begin{proof}
The case $R=0$ is immediate. For $|z|R\le1/2$, expansion of the
logarithm, with its linear term zero, gives
\[
 \left|\prod_i(1+zy_i)\right|\le \exp(NR^2|z|^2).
\]
Cauchy's coefficient bound on the circle
$|z|=\sqrt{k/(2N)}/R$ gives
$|e_k|\le(e^{1/2}R\sqrt{2N/k})^k$.
Divide by $\binom Nk\ge(N/k)^k$.
\end{proof}

Fix from now on
\[
 \ell=\lfloor m/4\rfloor,\quad N=m-\ell,\quad
 \alpha=\ell/m,\quad \theta=1/16.
\]
Let $B$ be a uniformly chosen $\ell$-element subset of $[m]$, and put
\[
 A_B=\prod_{j\in B}p_j,\qquad \overline A=\E_B A_B,
 \qquad x=\alpha s.
\]
The measure $\mu$ on rows defined by
\[
 \mu(q)=(\ell+1)w_q\overline A(q)
\]
is a probability measure. Define also the probability density
\[
 b_\ell(z)=\frac{\ell+1}{\ell}(1-z/\ell)^\ell
                 \boldsymbol1_{[0,\ell]}(z).
\]
Its expectation functional is denoted by $\beta_\ell$.

Let $L_j$ be the ordinary Laguerre polynomials, orthonormal for
$e^{-z}\,dz$ on $[0,\infty)$.

\begin{lemma}[Comparison with a scalar endpoint measure]\label{comparison}
Let $1\le d\le m^{1/4}$ and
\[
 Q(z)=\sum_{j=0}^d c_jL_j(\theta z),\qquad
 B_0=\sum_j|c_j|.
\]
For $P=Q^2$, and also for $P=(a-z)Q^2$ with $0\le a\le32$, one has
\[
 \left|\E_\mu P(x)-\beta_\ell(P)\right|
 \le C B_0^2\left(\frac d{\sqrt m}+\frac{d^2}{m}+e^{-cm}\right).
\tag{*}
\]
The constants are uniform in the coefficients and in $a$.
\end{lemma}
\begin{proof}
The classical bound $|L_j(z)|\le e^{z/2}$ for $z\ge0$ has a
self-contained proof in
\path{rational_designs/asymptotic/logarithmic_moment_barrier.tex}.
The derivative identity
\[
 L_j^{(h)}(z)=(-1)^hL_{j-h}^{(h)}(z),\qquad
 L_{j-h}^{(h)}=\sum_{v=0}^{j-h}\binom{j-v-1}{h-1}L_v
 \quad(h\ge1)
\]
therefore gives $|L_j^{(h)}(z)|\le\binom jh e^{z/2}$.
Consequently
\[
 |Q^{(h)}(z)|\le B_0\frac{d^h}{h!}e^{\theta z/2}.
\]
The product rule and
$\sum_{h=0}^k\binom kh^2=\binom{2k}k\le4^k$ imply
\begin{equation}\label{squarederivatives}
 \frac{|(Q^2)^{(k)}(z)|}{k!}
 \le B_0^2 e^{\theta z}\frac{(4d)^k}{(k!)^2}.
\end{equation}
For $P=(a-z)Q^2$, the corresponding upper bound is
\begin{equation}\label{linear-squarederivatives}
 B_0^2 e^{\theta z}\left[
  (a+z)\frac{(4d)^k}{(k!)^2}
  +\frac{(4d)^{k-1}}{((k-1)!)^2}\right],
\end{equation}
where the second term is omitted when $k=0$.

For fixed $B$, put $z_j=\ell u_j$ for $j\notin B$, and let
$v_B=N^{-1}\sum_{j\notin B}z_j$ be their mean.
Multilinearity and the exact mixed moments give
\begin{equation}\label{exactpolarization}
 (\ell+1)\E_w\E_B\left[A_B\Pol_N(P)(z_{[m]\setminus B})\right]
       =\beta_\ell(P).
\end{equation}
Indeed, after expanding the polynomial, every distinct-coordinate
product can be replaced by the corresponding product of one common
uniform variable $t$. The right side becomes
$(\ell+1)\int_0^1t^\ell P(\ell(1-t))\,dt$.
The degree condition holds because $\deg P\le2d+1\le N$ for large $m$.

On every row, \eqref{balance} gives
\[
 |z_j-v_B|\le C(s+1),\qquad
 0\le x,v_B\le\frac{\alpha}{1-\alpha}s.
\]
Also Maclaurin's inequality gives, on every row,
\begin{equation}\label{abar}
 \overline A=\frac{e_\ell(p)}{\binom m\ell}
        \le(1-s/m)^\ell\le e^{-\alpha s}.
\end{equation}
Taylor expansion of polarization about $v_B$ is exact:
\[
 \Pol_N(P)(z)-P(v_B)
  =\sum_{k=2}^{\deg P}\frac{P^{(k)}(v_B)}{k!}
                   \frac{e_k(z-v_B)}{\binom Nk}.
\]
Apply Lemma~\ref{elementary} and the derivative estimates.
Since $\theta\alpha/(1-\alpha)<\alpha/4$, the factor
$\E_B A_Be^{\theta v_B}$ is at most $e^{-3\alpha s/4}$.
Using \eqref{weightedmoments} and multiplying by $\ell+1$ bounds the
integrated error by a constant times
\[
 B_0^2\sum_{k=2}^{2d+1}
 \left[
  \frac{(Cd)^k(k+1)!}{(k!)^2}
  +\frac{(Cd)^{k-1}k!}{((k-1)!)^2}
 \right]\left(\frac{k}{m}\right)^{k/2}.
\]
The second term is needed only for a linear factor, but bounds both
cases. From $k!\ge(k/e)^k$ and $d\ge1$, the summand is at most
$C(k+1)^2(Cd/\sqrt m)^k$. For sufficiently large $m$ and
$d\le m^{1/4}$, this sum is at most $Cd^2/m$. Thus the polarization
part of the error is at most $CB_0^2d^2/m$.

It remains to replace $P(v_B)$ by $P(x)$. Under the uniform choice
of $B$, $\E_Bv_B=x$, and the variance formula for sampling without
replacement gives
\[
 \E_B(v_B-x)^2
  =\frac{\ell}{N(m-1)}\frac1m\sum_{j=1}^m(z_j-x)^2
  \le\frac{C(s+1)^2}{m},
\]
where in this formula $z_j=\ell u_j$ is defined for all $j$.
Because $0\le A_B\le1$, Cauchy--Schwarz implies
\[
 \E_B A_B|v_B-x|
 \le\sqrt{\overline A}\sqrt{\E_B(v_B-x)^2}
 \le e^{-\alpha s/2}\frac{C(s+1)}{\sqrt m}.
\]
The first-derivative cases of \eqref{squarederivatives} and
\eqref{linear-squarederivatives}, throughout the interval between
$x$ and $v_B$, are bounded by
\[
 CB_0^2d(1+s)e^{\theta\alpha s/(1-\alpha)}.
\]
The remaining exponential factor is at most $e^{-\alpha s/4}$.
After \eqref{weightedmoments} and multiplication by $\ell+1$,
this error is at most $CB_0^2d/\sqrt m$.

Together with \eqref{exactpolarization}, these estimates prove (*),
even without its harmless exponentially small term.
\end{proof}

\section*{Endpoint tests with an exponential margin}
For $d\ge1$ define
\[
 T_d(y)=\sum_{j=0}^{d-1}(1-j/d)L_j(y),\qquad
 R_d(z)=\left(\frac{2}{\theta d}-z\right)T_d(\theta z)^2.
\]
Laguerre orthogonality and the three-term recurrence give
\[
 \int_0^\infty e^{-y}T_d(y)^2\,dy
   =\frac{(d+1)(2d+1)}{6d},\qquad
 \int_0^\infty ye^{-y}T_d(y)^2\,dy
   =\frac{d+1}{2d}.
\]
Thus the mean of $z$ under $e^{-\theta z}T_d(\theta z)^2\,dz$
is $3/[\theta(2d+1)]\le3/(2\theta d)$.
The density ratio $b_\ell(z)/e^{-\theta z}$ is nonincreasing,
including its drop to zero beyond $\ell$. Multiplication by this
ratio cannot increase the mean: this follows directly by integrating
$(z-z')(r(z)-r(z'))\le0$ against two copies of the original measure.
Consequently
\[
 \beta_\ell(R_d)
 \ge\frac1{2\theta d}\beta_\ell(T_d(\theta z)^2).
\]
For $0\le z\le1/(4d)$ and $j<d$,
$|L_j(\theta z)-1|\le e^{j\theta z}-1<1/3$.
Hence $T_d(\theta z)\ge d/3$ there. Moreover $b_\ell(z)\ge3/4$
on that interval, by $(1-z/\ell)^\ell\ge1-z$.
It follows that
\begin{equation}\label{positiveR}
 \beta_\ell(R_d)\ge\frac1{96\theta}.
\end{equation}

Choose $d=\lfloor\kappa m^{1/6}\rfloor$, where $\kappa>0$ is a
sufficiently small absolute constant. The coefficient sum of $T_d$
is $(d+1)/2\le d$. Lemma~\ref{comparison} therefore bounds the error
for $R_d$ by
\[
 C\left(\frac{d^3}{\sqrt m}+\frac{d^4}{m}+d^2e^{-cm}\right).
\]
By first choosing $\kappa$ small and then $m$ large, this is smaller
than half the positive constant in \eqref{positiveR}.
Thus $\E_\mu R_d(x)>0$. Since $R_d$ is nonpositive for
$x\ge2/(\theta d)$, some row of positive $\mu$-mass lies below
that threshold. The last row minimizes $s$ and $x$, so
\begin{equation}\label{lastx}
 x_K<2/(\theta d).
\end{equation}

Put $r=\lfloor d/8\rfloor$ and
\[
 S=\sum_{j=0}^{r}L_j(\theta x_K)^2,
 \qquad Q_*(z)=\frac1S\sum_{j=0}^{r}
              L_j(\theta x_K)L_j(\theta z).
\]
Then $Q_*(x_K)=1$. From \eqref{lastx} we have
$j\theta x_K\le1/4$, so $|L_j(\theta x_K)-1|<1/3$.
In particular $S\ge c d$, and the sum of absolute Laguerre
coefficients of $Q_*$ is bounded by an absolute constant.
Since $b_\ell(z)\le2e^{-z}\le2e^{-\theta z}$,
orthogonality gives
\[
 \beta_\ell(Q_*^2)\le\frac2{\theta S}\le C/d.
\]
Lemma~\ref{comparison}, uniformly also for this data-dependent
polynomial, implies
\[
 \mu(K)\le\E_\mu Q_*(x)^2
 \le \frac Cd+C\left(\frac d{\sqrt m}+\frac{d^2}{m}+e^{-cm}\right)
 \le C/d.
\]
Finally $s_K=x_K/\alpha=O(1/d)$. For every $B$,
$A_B(K)=\prod_{j\in B}(1-u_j(K))\ge1-s_K\ge1/2$ for large $m$.
Hence
\[
 (\ell+1)w_K/2\le\mu(K)\le C/d,
 \qquad w_K\le C m^{-7/6}.
\]
Reversing the rows and replacing every $p_j$ by $1-p_j$ gives the
same bound for $w_1$. For a distinguished uniform die, $w_q=1/K$;
the standard comparison tensor has $m=n-1$ and the stated mixed moments.
This proves Theorem~\ref{main}, as claimed.
\end{document}
